%% ma: L10-B12-0001
%% lop: 10
%% chuong: 5
%% bai: 12
%% dang: 10.12.1
%% loai: TL
%% muc: NB
%% dap_an: "-"
%% nguon: "chuyen-de-cac-so-dac-trung-cua-mau-so-lieu-khong-ghep-nhom-toan-10-chuong-trinh-moi.pdf (GV. Phan Nhật Linh) (P01), Bài 01 Bài tập tự luận 1"
%% kien_thuc: "Lớp 10 Bài 12 (số gần đúng, sai số tuyệt đối, độ chính xác, sai số tương đối, quy tròn số)"
%% trang_thai: chinh-thuc
%% ghi_chu: "Tự duyệt theo yêu cầu giáo viên 22/10/2026."
\begin{ex}
Theo em, các con số sau đây là số đúng hay số gần đúng?

a) Theo số liệu của Liên Hợp Quốc ngày $08/8/2022$, dân số của Việt Nam là $99\,032\,076$ người.

b) Hóa đơn tiền điện tháng $8$ năm $2022$ của gia đình bác Hoa là $525\,314$ đồng.

c) Biết $\sqrt[3]{7}=1{,}9129311\ldots$ Viết số gần đúng của $\sqrt[3]{7}$ khi làm tròn đến hai chữ số thập phân, ba chữ số thập phân.

d) Một thửa đất hình vuông có cạnh $135$ m. Biết $\sqrt2=1{,}4142135\ldots$ Tính độ dài đường chéo của thửa đất (làm tròn đến hàng chục).
\loigiai{a) Số gần đúng (dân số thay đổi liên tục, chỉ ước tính được).

b) Số đúng.

c) $\sqrt[3]7\approx1{,}91$ (hai chữ số thập phân); $\sqrt[3]7\approx1{,}913$ (ba chữ số thập phân).

d) Đường chéo $135\sqrt2\approx190{,}918\ldots$; làm tròn đến hàng chục được $190$ m.}
\end{ex}
